\section{구조 정리: 중복 제거, 블록 방식, greedy 사다리, 매끄러운 수 경로 (STRUCT)}
\label{STRUCT:sec}

이 절은 Round 4 설계의 S1, S2, S3, S6을 처음부터 다시 증명하고, 어떤 산술적 명제가 어떤 상계를 주는지를 정확한 함의로 정리한다.
Erdős--Graham \#304 자체는 \OPEN 이며 여기서 해결을 주장하지 않는다.
모든 명제에는 꼬리표가 붙어 있다. 수치 확인은 \texttt{\$R/agents/STRUCT/code/}의 스크립트와 \texttt{\$R/agents/STRUCT/out/}의 출력에 있다
(\texttt{\$R} $=$ \texttt{/home/user/FORMALIZATION-OF-MY-10-PAPERS/erdos304/round4}).

\paragraph{기호.}
$N\ge 2$에 대해 $D(N)$은 $N$의 양의 약수 전체, $D^*(N)=D(N)\setminus\{N\}$이다.
정수 $a\ge 0$에 대해
\[
g_N(a)=\min\Bigl\{|S| : S\subseteq D^*(N),\ \textstyle\sum_{d\in S}d=a\Bigr\},\qquad
\widetilde g_N(a)=\min\Bigl\{s : a=d_1+\dots+d_s,\ d_i\in D^*(N)\Bigr\}
\]
($\min\emptyset=\infty$, $g_N(0)=\widetilde g_N(0)=0$)로 둔다. $\widetilde g_N$에서는 같은 약수의 반복을 허용한다.
$\Lam_L=\lcm(1,\dots,L)$, $a_p=\lfloor \log L/\log p\rfloor$, $\psi(L)=\log\Lam_L$, $g_L=g_{\Lam_L}$, $D^*_L=D^*(\Lam_L)$,
$H(L)=\max_{0\le a<\Lam_L}g_L(a)$, $P=\{p:\ L/2<p\le L\}$, $M=|P|$이다.
$e(x)=\exp(2\pi i x)$이다. 집합 $A\subseteq\Z_{\ge0}$와 $k\ge 0$에 대해 $kA=\{x_1+\dots+x_k: x_i\in A\}$ ($0A=\{0\}$)는 반복을 허용한 합집합(sumset)이다.
외부 입력으로 쓰는 소수정리(PNT)는 다음 형태이다.

\begin{theorem}[소수정리, 교과서]\label{STRUCT:ext-pnt}
\EXT\ (Montgomery--Vaughan, \emph{Multiplicative Number Theory I}, 6장; 교과서 사실, 정리 번호는 recalled)
$y\to\infty$일 때 $\psi(y)=(1+o(1))y$, $\theta(y)=\sum_{p\le y}\log p=(1+o(1))y$, $\pi(y)=(1+o(1))y/\log y$.
또 모든 $y\ge 2$에서 $0\le\psi(y)-\theta(y)=\sum_{k\ge2}\theta(y^{1/k})\le \sqrt y\,(\log y)^2$ (초등적: Chebyshev의 $\theta(x)<x\log 4$와 $k\le\log_2 y$에서).
\end{theorem}

%======================================================================
\subsection{S1: 최소공배수형 중복 제거와 다중집합 정식화}
\label{STRUCT:sec-S1}

\begin{definition}[성질 $(\star)$]\label{STRUCT:def-star}
정수 $N\ge2$가 성질 $(\star)$를 가진다는 것은 $N$을 나누는 모든 홀수 소수 $p$에 대해 $r_p:=(p+1)/2$가 $N$을 나눈다는 뜻이다.
\end{definition}

\begin{lemma}\label{STRUCT:lem-star-examples}
\PROVED\ 모든 $L\ge 2$에 대해 $\Lam_L$은 $(\star)$를 가진다. 모든 $n\ge2$에 대해 $n!$과 $2^n$도 $(\star)$를 가진다. $N=10$은 $(\star)$를 가지지 않는다.
\end{lemma}
\begin{proof}
홀수 소수 $p\mid\Lam_L$이면 $p\le L$이므로 $r_p=(p+1)/2\le L$이고, $1,\dots,L$은 모두 $\Lam_L$을 나누므로 $r_p\mid\Lam_L$.
$n!$도 같다($p\le n\Rightarrow r_p\le n\Rightarrow r_p\mid n!$). $2^n$에는 홀수 소인수가 없다.
$N=10$에서 $p=5$이면 $r_5=3\nmid 10$.
\end{proof}

\begin{lemma}[중복쌍 교환]\label{STRUCT:lem-swap}
\PROVED\ $N$이 $(\star)$를 가지고 $d\in D(N)$, $2d<N$이라 하자. $p:=\min\{q \text{ 소수}: v_q(d)<v_q(N)\}$로 두면 $p$는 존재하고 $p\mid N$이다.
\begin{enumerate}[label=(\roman*)]
\item $p=2$이면 $2d\in D^*(N)$.
\item $p\ge3$이면 $r=(p+1)/2$는 $d$를 나누고, $u:=pd/r$, $v:=d/r$는 $D^*(N)$의 원소이며
\[
u+v=2d,\qquad 0<v<d<u<2d,\qquad u^2+v^2-2d^2=2(d-v)^2\ge2,\qquad \frac{uv}{d^2}=\frac{4p}{(p+1)^2}\le\frac34 .
\]
\end{enumerate}
\end{lemma}
\begin{proof}
$d\mid N$이므로 모든 소수 $q$에서 $v_q(d)\le v_q(N)$이다. 모든 $q$에서 등호이면 $d=N$인데 $2d<N$이므로 불가능하다. 따라서 $p$가 존재하고, $v_p(N)>v_p(d)\ge0$이므로 $p\mid N$이다.

(i) $v_2(2d)=v_2(d)+1\le v_2(N)$이고 홀수 소수 $q$에서는 $v_q(2d)=v_q(d)\le v_q(N)$이므로 $2d\mid N$. $2d<N$이므로 진약수이다.

(ii) $p$는 $N$을 나누는 홀수 소수이므로 $(\star)$에 의해 $r\mid N$. $p\ge3$이면 $2\le r<p$이다. $r$의 임의의 소인수 $q$는 $q\le r<p$이므로 $p$의 최소성에 의해 $v_q(d)=v_q(N)$이고, $r\mid N$에서 $v_q(r)\le v_q(N)=v_q(d)$. 따라서 $r\mid d$이고 $u,v$는 정수이다.
$v\mid d\mid N$. $r<p$이므로 $v_p(r)=0$이고 $v_p(u)=v_p(d)+1\le v_p(N)$; $q\ne p$에서는 $v_q(u)=v_q(d)-v_q(r)\le v_q(N)$. 그러므로 $u\mid N$.
$u+v=(p+1)d/r=2d$. $r\ge2$에서 $v\le d/2<d$, 따라서 $u=2d-v>d$이고 $u<2d<N$. 그러므로 $u,v\in D^*(N)$, $u\ne v$.
$u^2+v^2-2d^2=(2d-v)^2+v^2-2d^2=2(d-v)^2$이고 $d-v$는 양의 정수이므로 $\ge2$.
끝으로 $uv/d^2=p/r^2=4p/(p+1)^2$이고 $3(p+1)^2-16p=(3p-1)(p-3)\ge0$ ($p\ge3$).
\end{proof}

\begin{theorem}[중복 제거; baker Thm~3.2의 일반화]\label{STRUCT:thm-dedup}
\PROVED\ $N$이 $(\star)$를 가지고 $0\le a<N$, $a=d_1+\dots+d_s$ ($d_i\in D^*(N)$, 반복 허용)라 하자. 그러면 $a$는 $D^*(N)$의 서로 다른 원소 많아야 $s$개의 합이다. 따라서 $g_N(a)=\widetilde g_N(a)$이다.
더 정확히, 아래의 교환 알고리즘은 많아야
\[
\min\Bigl\{\tfrac{a^2-\Phi_0}{2},\ (s-1)+\tfrac{\sum_i\log d_i+(s-1)\log2}{\log(4/3)}\Bigr\},\qquad \Phi_0=\sum_i d_i^2,
\]
번의 교환 후 종료하며 그 출력이 원하는 구별 표현이다.
\end{theorem}
\begin{proof}
$a=0$이면 빈 표현이다. $a\ge1$이라 하자. 상태는 $D^*(N)$의 양의 원소들로 이루어진 유한 다중집합 $\mathcal M$으로, 불변식 (I1) 원소의 합 $=a$, (I2) $|\mathcal M|\le s$를 만족한다. 초기 상태는 $\{d_1,\dots,d_s\}$이다.

\emph{교환.} 어떤 $d$가 $\mathcal M$에서 중복도 $\ge2$를 가지면, 모든 원소가 양수이므로 $2d\le a<N$이다. 보조정리~\ref{STRUCT:lem-swap}의 $p$를 잡아 $p=2$이면 두 개의 $d$를 하나의 $2d$로(병합), $p\ge3$이면 두 개의 $d$를 $u,v$로(분할) 바꾼다. 새 원소들은 $D^*(N)$의 양의 원소이고 합은 그대로이며, 원소 수는 병합에서 1 감소, 분할에서 불변이다. 따라서 (I1), (I2)가 유지된다.

\emph{종료 (제곱합).} $\Phi(\mathcal M)=\sum_{x\in\mathcal M}x^2$ (중복도 포함)는 정수이고, 양수들에 대해 $\sum x^2\le(\sum x)^2$이므로 항상 $\Phi\le a^2$. 병합은 $\Phi$를 $(2d)^2-2d^2=2d^2\ge2$만큼, 분할은 보조정리~\ref{STRUCT:lem-swap}(ii)에 의해 $\ge2$만큼 늘린다. 따라서 교환 횟수는 $(a^2-\Phi_0)/2$ 이하이다.

\emph{종료 (곱).} $\Pi(\mathcal M)=\prod_{x\in\mathcal M}x\ge1$ (양의 정수들의 곱). 병합은 $\Pi$에 $2d/d^2=2/d\le2$를, 분할은 $uv/d^2\le3/4$를 곱한다. 병합은 원소 수를 줄이고 원소 수는 $\ge1$로 남으므로 (합 $a\ge1$) 병합 횟수 $M_{\rm g}\le s-1$. 임의의 실행 접두에서 분할 횟수를 $R$이라 하면 $1\le\Pi\le\Pi_0\,2^{M_{\rm g}}(3/4)^R$이므로 $R\le(\log\Pi_0+(s-1)\log2)/\log(4/3)$. 이 부등식은 종료를 가정하지 않고 모든 접두에서 성립하므로 두 번째 종료 증명이 된다.

\emph{출력.} 종료 상태에서는 중복도 $\ge2$인 원소가 없다(있으면 교환이 가능하다). 따라서 종료 상태는 $D^*(N)$의 서로 다른 원소 $\le s$개로 $a$를 표현한다. 이로써 $g_N(a)\le\widetilde g_N(a)$이고, 역방향 $\widetilde g_N\le g_N$은 구별 표현이 곧 반복 표현이므로 자명하다.
\end{proof}

\begin{theorem}[$(\star)$에 의한 완전한 특징짓기]\label{STRUCT:thm-char}
\PROVED\ $N\ge2$에 대해 다음은 동치이다.
(i) 모든 $0\le a<N$에서 $g_N(a)=\widetilde g_N(a)$;
(ii) $N$은 $(\star)$를 가진다.
더구나 $(\star)$가 홀수 소수 $p\mid N$에서 깨지면 ($r_p\nmid N$), $a^*:=2N/p$는 $\widetilde g_N(a^*)=2<3\le g_N(a^*)$ (값 $\infty$ 포함)를 만족한다.
\end{theorem}
\begin{proof}
(ii)$\Rightarrow$(i)는 정리~\ref{STRUCT:thm-dedup}. (i)$\Rightarrow$(ii)의 대우를 보인다. 홀수 소수 $p\mid N$에서 $r=r_p\nmid N$이라 하자.
$a^*=2N/p$는 정수이고 $p\ge3$이므로 $a^*<N$. $N/p\in D^*(N)$이고 $a^*=N/p+N/p$이므로 $\widetilde g_N(a^*)\le2$. $N/a^*=p/2\notin\Z$이므로 $a^*$는 약수가 아니고, $a^*\ne0$이므로 $\widetilde g_N(a^*)=2$.
이제 $a^*=e+f$, $e<f$, $e,f\in D(N)$이라 가정하자. $f>a^*/2=N/p$이고 $f<a^*=2N/p$. $m:=N/f$는 $N$의 약수인 정수이고 $p/2<m<p$이다. $p$가 홀수이므로 $2m\ge p+1$, 즉 $j:=2m-p$는 $1\le j\le p-2$인 정수이다.
$e=2N/p-N/m=N j/(pm)$이고 $e\mid N$이므로 $N/e=pm/j$는 정수이다. $1\le j<p$이고 $p$가 소수이므로 $\gcd(j,p)=1$, 따라서 $j\mid m$. 그러면 $j\mid 2m-j=p$이므로 $j\in\{1,p\}$이고 $j<p$에서 $j=1$, 즉 $m=(p+1)/2=r$. 그러나 $m\mid N$이므로 $r\mid N$, 모순. 따라서 $a^*$는 서로 다른 약수 두 개 이하의 합이 아니고 $g_N(a^*)\ge3$.
\end{proof}

\begin{counterexample}[$N=10$, $a=4$]\label{STRUCT:cex-10}
\REFUTED\ ``모든 $N$에서 반복 허용 최소 항수 $=$ 구별 최소 항수''는 거짓이다.
$D^*(10)=\{1,2,5\}$의 구별 부분합 전체는 $\{0,1,2,3,5,6,7,8\}$이므로 $g_{10}(4)=g_{10}(9)=\infty$이지만 $4=2+2$, $9=5+2+2$로 $\widetilde g_{10}(4)=2$, $\widetilde g_{10}(9)=3$.
증명의 실패 지점: $d=2$의 최소 부족 소수는 $p=5$이고 $r=3\nmid 10$이다(정리~\ref{STRUCT:thm-char}의 $a^*=2\cdot10/5=4$).
\end{counterexample}

\begin{remark}[수치 확인]\label{STRUCT:rem-num-dedup}
\NUMERIC\ \texttt{code/dedup\_check.py} (출력 \texttt{out/dedup\_check.out}):
(a) $2\le L\le16$의 모든 $0\le a<\Lam_L$에서 0/1 배낭 DP로 계산한 $g_L(a)$와 sumset BFS로 계산한 $\widetilde g_L(a)$가 일치하며 $H(L)=1,2,3,4,4,5,5,5,5,6,6,6,6,6,6$;
(b) $2\le N\le6000$에서 (i)$\Leftrightarrow$(ii)가 예외 없이 성립($(\star)$를 가지는 $N$ 283개, 가지지 않는 $N$ 5716개), 그리고 $(\star)$가 깨지는 9286개의 쌍 $(N,p)$ 모두에서 $a^*=2N/p$가 $\widetilde g=2$, $g\ge3$;
(c) $L\le60$의 무작위 반복 표현 3000개에서 교환 알고리즘이 합·구별성·항수 비증가를 지키고 두 교환 횟수 상계를 만족(분할 수/곱 상계의 최댓값 0.138);
baker 예 $N=60$: $12+12+4+20\mapsto 2+6+10+30$.
\end{remark}

\begin{corollary}[$(\star)$이면 실용수]\label{STRUCT:cor-practical}
\PROVED\ $N$이 $(\star)$를 가지면 모든 $1\le a<N$은 $N$의 서로 다른 약수의 합이다. 특히 $\Lam_L$, $n!$은 이 의미에서 실용수이다(p1p4 Lemma~8.3의 세 번째 증명).
\end{corollary}
\begin{proof}
$a=1+\dots+1$ ($a$개)은 반복 표현이고 $1\in D^*(N)$. 정리~\ref{STRUCT:thm-dedup}.
\end{proof}

\begin{corollary}[다중집합·sumset 정식화]\label{STRUCT:cor-sumset}
\PROVED\ $A_L:=D^*_L\cup\{0\}$라 하면 $0\le a<\Lam_L$, $k\ge0$에서 $g_L(a)\le k\iff a\in kA_L$. 따라서
\[
H(L)=\min\{k\ge0:\ kA_L\supseteq[0,\Lam_L)\cap\Z\}.
\]
또한 (a) $x+y<\Lam_L$이면 $g_L(x+y)\le g_L(x)+g_L(y)$; (b) Bellman 점화식 $g_L(a)=1+\min_{d\in D^*_L,\,d\le a}g_L(a-d)$ ($1\le a<\Lam_L$);
(c) $[z^a]\bigl(1+\sum_{d\in D^*_L}z^d\bigr)^k>0\iff g_L(a)\le k$ ($0\le a<\Lam_L$).
\end{corollary}
\begin{proof}
$g_L(a)\le k$이면 구별 표현에 $0$을 채워 $a\in kA_L$. 역으로 $a\in kA_L$이면 $0$을 지운 것은 $\le k$항의 반복 표현이므로 $\widetilde g_L(a)\le k$이고 정리~\ref{STRUCT:thm-dedup}와 보조정리~\ref{STRUCT:lem-star-examples}로 $g_L(a)\le k$.
(a) 두 최적 표현을 합친 반복 표현에 정리~\ref{STRUCT:thm-dedup}. (b) 최적 표현에서 한 항 $d$를 빼면 $g_L(a-d)\le g_L(a)-1$; 역으로 $a-d$의 최적 표현에 $d$를 붙인 반복 표현에 정리~\ref{STRUCT:thm-dedup}. (c) 좌변은 $k$개 자리에 $0$ 또는 약수를 넣는 순서 있는 선택의 수이므로 $a\in kA_L$과 동치.
\end{proof}

\begin{theorem}[서로소성 없는 cascade 조립]\label{STRUCT:thm-cascade}
\PROVED\ $N$이 $(\star)$를 가진다고 하자. $I_0,\dots,I_{K+1}\subseteq\Z_{\ge0}$, $I_0\subseteq[0,N)$, $I_{K+1}=\{0\}$, 임의의 (서로 겹쳐도, 심지어 같아도 되는) 족 $\mathcal D_0,\dots,\mathcal D_K\subseteq D^*(N)$, 예산 $s_k\ge0$이 다음을 만족한다고 하자:
\[
(\mathrm{P1}_k^{\rm rep})\quad \forall R\in I_k\ \exists\ \text{다중집합 } S\subseteq\mathcal D_k,\ |S|\le s_k,\ R-\textstyle\sum S\in I_{k+1}.
\]
(선택 $S$는 이전 단계의 선택, 곧 이력에 의존해도 된다.) 그러면 모든 $a\in I_0$에서 $g_N(a)\le\sum_k s_k$.
\end{theorem}
\begin{proof}
$R_0=a$에서 시작하여 $R_{k+1}=R_k-\sum S_k$로 진행하면 $R_{K+1}=0$이고 $a=\sum_k\sum S_k$는 $D^*(N)$의 원소 $\le\sum_k s_k$개의 반복 표현이다. $a<N$이므로 정리~\ref{STRUCT:thm-dedup}.
\end{proof}

\begin{remark}[p1p4 반례 8.12와 설계 조건의 제거]\label{STRUCT:rem-8-12}
\PROVED\ p1p4 반례~8.12 ($n=6$, $I_0=\{4\}$, $I_1=\{2\}$, $\mathcal D_0=\mathcal D_1=\{2\}$)는 ``이어 붙인 출력 자체가 구별 표현''이라는 문장만 반박한다. $6=\Lam_3$은 $(\star)$를 가지므로 결론 $g_6(4)\le2$는 참이다: $d=2$에서 $p=3$, $r=2$, $4=3+1$.
반면 $(\star)$가 없는 $N=10$에서 같은 자료 ($I_0=\{4\}$, $I_1=\{2\}$, $I_2=\{0\}$, $\mathcal D_0=\mathcal D_1=\{2\}$, $s_0=s_1=1$)는 단계성질을 모두 만족하지만 $g_{10}(4)=\infty$이다. 즉 $(\star)$는 정리~\ref{STRUCT:thm-cascade}에 본질적이다.
따라서 $N=\Lam_L$에서는 p1p4의 (SP1-disj) (가설 8.19), 라벨 서로소 설계(명제 8.21), 이력의존 회피(정리 8.11), 반복 호출의 $u$-강건성(따름정리 8.15(b))이 모두 불필요하다: p1p4 정리 8.20은 (SP1-top), (SP1-LOS), (SP1-small)만으로 성립하고, 같은 메커니즘의 $J$번 호출은 강건성 조건 없이 예산 $J s$로 합산된다.
\end{remark}

\begin{corollary}[순서 있는 튜플 계수로 충분; Newton 보정 불필요]\label{STRUCT:cor-newton}
\PROVED\ $N$이 $(\star)$를 가지고 $\mathcal D\subseteq D^*(N)$, $T(\theta)=\sum_{d\in\mathcal D}e(d\theta)$, $r_K(n)=\int_0^1T(\theta)^Ke(-n\theta)\,d\theta=\#\{(d_1,\dots,d_K)\in\mathcal D^K:\sum d_i=n\}$라 하자. $0\le n<N$이고 $r_K(n)>0$이면 $g_N(n)\le K$.
따라서 baker 12.3절의 Newton 항등식 $24E_4=T^4-6T^2T(2\theta)+3T(2\theta)^2+8TT(3\theta)-6T(4\theta)$와 그 보정항 추정은 불필요하며, 블록 응용(baker 정리 13.1)에서는 금지집합 $U$를 피할 필요도 없다(블록 사이의 겹침은 정리~\ref{STRUCT:thm-cascade}가 처리한다).
\end{corollary}
\begin{proof}
순서 있는 튜플은 반복 표현이다. 정리~\ref{STRUCT:thm-dedup}.
\end{proof}

\begin{remark}[중복 제거가 보존하지 않는 것]\label{STRUCT:rem-not-preserved}
\PROVED\ 교환의 새 항 $u=pd/r$, $v=d/r$는 미리 정한 족 $\mathcal D$ 밖에 있을 수 있다. 따라서 ``$\mathcal D$의 원소만 쓰는 구별 표현''(크기창, 금지집합 $U$ 회피, 풀 소수 개수 $\omega_P$ 고정 등)은 반복 표현으로부터 따라 나오지 않는다. 예: $N=6$, $\mathcal D=\{2\}$에서 $4=2+2$이지만 $4$는 $\mathcal D$의 구별 합이 아니다.
또 정리~\ref{STRUCT:thm-dedup}는 항 수를 늘리지 않을 뿐, 이미 서로 다른 $k+1$개의 항을 $k$개로 줄이지 않는다. 제어되는 것은 제한 없는 $g_N$뿐이다.
\end{remark}

%======================================================================
\subsection{S2: 블록(묶음 혼합진법) 방식}
\label{STRUCT:sec-S2}

\begin{definition}[블록 방식]\label{STRUCT:def-block}
$\Lam=\Lam_L$에 대한 블록 방식은 약수 사슬 $1=N_0\mid N_1\mid\dots\mid N_t=\Lam$ ($t\ge1$, $b_j:=N_j/N_{j-1}\ge2$)과 예산 $K_1,\dots,K_t\in\Z_{\ge0}$로서, $W_j:=\Lam/N_j$라 할 때 다음 단계성질을 만족하는 것이다.
\[
(\mathrm B_j)\qquad \text{모든 정수 } 0\le c<b_j \text{는 } D(N_j) \text{의 원소 많아야 } K_j \text{개의 합이다 (반복 허용).}
\]
비용은 $\sum_jK_j$이다. 모든 $j$에서 $N_j=\Lam_{y_j}$ ($y_1<\dots<y_t=L$, $\Lam_{y_1}<\dots<\Lam_{y_t}$)이면 $\Lam$형이라 부르고
\[
\nu_j:=\frac{\log b_j}{\psi(y_j)}\in(0,1],\qquad s_j:=\log\frac{\psi(y_j)}{\psi(y_{j-1})}=\log\frac1{1-\nu_j}\quad(j\ge2)
\]
로 둔다($\nu_1=1$). $y'\le y$에 대해 $T(y,y'):=\max_{0\le c<\Lam_y/\Lam_{y'}}g_y(c)$ ($T(y,1)=H(y)$)로 둔다.
\end{definition}

\begin{theorem}[블록 방식 상계]\label{STRUCT:thm-block}
\PROVED\ 블록 방식이 있으면 $H(L)\le\sum_{j=1}^tK_j$.
\end{theorem}
\begin{proof}
$W_0=\Lam$, $W_t=1$, $W_{j-1}=b_jW_j$. $0\le a<\Lam$에 대해 $a_0=a$, $c_j=\lfloor a_{j-1}/W_j\rfloor$, $a_j=a_{j-1}-c_jW_j$로 두면 $0\le a_j<W_j$이고 $a_{j-1}<W_{j-1}=b_jW_j$에서 $0\le c_j<b_j$; $a_t<W_t=1$이므로 $a=\sum_jc_jW_j$.
$(\mathrm B_j)$로 $c_j=\sum_ie_{j,i}$ ($\le K_j$항, $e_{j,i}\in D(N_j)$, $e_{j,i}\le c_j$). $e_{j,i}W_j\mid N_jW_j=\Lam$이고 $e_{j,i}W_j\le c_jW_j\le a<\Lam$이므로 $e_{j,i}W_j\in D^*_L$. 따라서 $a$는 $D^*_L$의 원소 $\le\sum K_j$개의 반복 표현이고, 정리~\ref{STRUCT:thm-dedup}와 보조정리~\ref{STRUCT:lem-star-examples}로 $g_L(a)\le\sum K_j$.
\end{proof}

\begin{corollary}[자기닮음 부등식]\label{STRUCT:cor-selfsim}
\PROVED\ $1\le y<L$이면 $H(L)\le H(y)+T(L,y)$이고 $T(L,y)\le H(L)$. $\Lam$형 방식에서 $(\mathrm B_j)$를 만족하는 최소 예산은 정확히 $K_1=H(y_1)$, $K_j=T(y_j,y_{j-1})$ ($j\ge2$)이며, 따라서 모든 사슬 $y_1<\dots<y_t=L$에 대해 $H(L)\le H(y_1)+\sum_{j\ge2}T(y_j,y_{j-1})$.
\end{corollary}
\begin{proof}
사슬 $1\mid\Lam_y\mid\Lam_L$에 정리~\ref{STRUCT:thm-block}: $c<\Lam_y$는 $\Lam_y$의 약수 $\le H(y)$개의 합이고, $c<\Lam_L/\Lam_y$는 정의상 $\Lam_L$의 약수 $\le T(L,y)$개의 합이다. $T\le H$는 $\Lam_L/\Lam_y\le\Lam_L$에서. 최소 예산: $c<b_j\le\Lam_{y_j}$이므로 $(\mathrm B_j)\iff K_j\ge\max_{c<b_j}\widetilde g_{\Lam_{y_j}}(c)$이고, $\Lam_{y_j}$에서 정리~\ref{STRUCT:thm-dedup}로 $\widetilde g=g$.
\end{proof}

\begin{proposition}[블록별 counting 필요조건]\label{STRUCT:prop-block-count}
\PROVED\ $(\mathrm B_j)$가 성립하면 $D_j:=\#\{e\in D(N_j):1\le e<b_j\}$에 대해 $\binom{D_j+K_j}{K_j}\ge b_j$, 특히 $K_j\ge\log b_j/\log(1+D_j)$.
\end{proposition}
\begin{proof}
각 $c<b_j$는 $\{0\}\cup\{e\mid N_j:1\le e<b_j\}$ ($D_j+1$개 원소)에서 뽑은 크기 정확히 $K_j$인 다중집합의 합이다(항은 $\le c<b_j$, 빈 자리는 $0$). 그런 다중집합의 수는 $\binom{D_j+K_j}{K_j}$이다. 또 $\binom{D+K}{K}=\prod_{i=1}^K\frac{D+i}{i}\le(1+D)^K$ ($D+i\le i(1+D)$).
\end{proof}

\begin{lemma}[약수 개수 상계; p1p4 보조정리 8.34의 재검증과 일반화]\label{STRUCT:lem-divcount}
\PROVED\ $y\ge16$, $0<\eta\le1/2$, $u>0$이고 $m_0:=u/((1-\eta)\log y)\le\pi(y)$이면
\[
\log\#\{e\mid\Lam_y: e\le \e^u\}\le m_0\log\frac{\e\,\pi(y)}{m_0}+\pi(y^{1-\eta})\log(1+\log_2y).
\]
\end{lemma}
\begin{proof}
$e\mid\Lam_y$를 $e=e_1e_2$로 쓴다: $e_1$은 $p\le y^{1-\eta}$ 부분, $e_2$는 $(y^{1-\eta},y]$의 소수 부분. $\eta\le1/2$이면 $p>y^{1-\eta}\ge\sqrt y$에서 $p^2>y$, 즉 $a_p=1$이므로 $e_2$는 서로 다른 소수들의 곱이다.
$e_1$의 선택 수는 $\prod_{p\le y^{1-\eta}}(a_p+1)\le(1+\log_2y)^{\pi(y^{1-\eta})}$ ($a_p\le\log_2y$).
$e_2$가 $m\ge1$개의 소수를 가지면 $e_2>y^{(1-\eta)m}$이고 $e_2\le e\le\e^u$이므로 $m<m_0$. 따라서 $e_2$의 선택 수는 $\sum_{0\le m\le m_0}\binom{\pi(y)}{m}$ 이하이다.
$0<k\le n$ (실수 $k$)에 대해 $\sum_{m\le k}\binom nm\le(\e n/k)^k$: $m\le k$에서 $(k/n)^k\le(k/n)^m$이므로 $(k/n)^k\sum_{m\le k}\binom nm\le\sum_m\binom nm(k/n)^m=(1+k/n)^n\le\e^k$. $n=\pi(y)$, $k=m_0$로 쓰면 결론.
\end{proof}

\begin{lemma}[$\Lam$형 블록 하나의 비용]\label{STRUCT:lem-perblock}
\PROVED\ (정리~\ref{STRUCT:ext-pnt} 사용.) $0<w<1$, $\eps>0$을 고정하고
\[
G(w):=\frac{1}{(1+\log\frac1w)\log\frac1{1-w}}
\]
라 하자. $Y_1=Y_1(w,\eps)$가 있어, $\Lam$형 방식의 블록 $j\ge2$가 $y_j\ge Y_1$, $\nu_j\le w$를 만족하면 $K_j\ge(1-\eps)\,G(w)\,s_j\log y_j$.
\end{lemma}
\begin{proof}
$y=y_j$, $\nu=\nu_j$, $\eta=\eta(y):=6(\log\log y)^2/\log y$, $\nu_*(y):=y^{-\eta/2}=\exp(-3(\log\log y)^2)$, $\rho(y):=\pi(y)\log y/\psi(y)$로 둔다. PNT로 $\eta\to0$, $\rho(y)\to1$, $\theta(y)/\psi(y)\to1$, $\psi(y)\ge y/2$ (큰 $y$), 그리고 $E(y):=\pi(y^{1-\eta})\log(1+\log_2y)\le y^{1-\eta}$ (큰 $y$; $\pi(z)\le2z/\log z$).

\emph{경우 A: $\nu\ge\nu_*(y)$.} $b_j\ge2$이므로 $D_j\ge1$이고 $\log(1+D_j)\le\log2+\log D_j$. $u=\log b_j=\nu\psi(y)$로 보조정리~\ref{STRUCT:lem-divcount}을 쓴다: $m_0=\nu\psi(y)/((1-\eta)\log y)\le w\psi(y)/((1-\eta)\log y)\le\pi(y)$는 $w\psi(y)\le(1-\eta)\theta(y)\le(1-\eta)\pi(y)\log y$에서 따르고, 이는 $w<1$과 $\theta/\psi\to1$, $\eta\to0$으로 큰 $y$에서 성립한다. $\e\pi(y)/m_0=\e(1-\eta)\rho(y)/\nu\le\e\rho(y)/\nu$이므로
\[
\log(1+D_j)\le\frac{\nu\psi(y)}{(1-\eta)\log y}\Bigl[\log\frac1\nu+\kappa\Bigr],\qquad
\kappa=\kappa(y):=1+|\log\rho(y)|+\frac{(\log2+E(y))\log y}{\nu_*(y)\psi(y)} .
\]
마지막 항은 $\le 2y^{1-\eta/2}\log y/(y/2)=4\log y\,\exp(-3(\log\log y)^2)\to0$이므로 $\kappa(y)\to1$이고 $\kappa\ge1$. 명제~\ref{STRUCT:prop-block-count}로
\[
K_j\ge\frac{\log b_j}{\log(1+D_j)}\ge\frac{(1-\eta)\log y}{\log(1/\nu)+\kappa}=(1-\eta)\,s_j\log y\;G_\kappa(\nu),\qquad G_\kappa(\nu):=\frac1{(\log\frac1\nu+\kappa)\log\frac1{1-\nu}} .
\]
$\kappa\ge1$이면 $G_\kappa$는 $(0,1)$에서 감소한다: 분모 $h(\nu)$의 도함수는 $h'(\nu)=-\frac1\nu\log\frac1{1-\nu}+\frac{\log(1/\nu)+\kappa}{1-\nu}$이고 $\log\frac1{1-\nu}=\sum_k\nu^k/k\le\nu/(1-\nu)$이므로 $\frac1\nu\log\frac1{1-\nu}\le\frac1{1-\nu}<\frac{\log(1/\nu)+\kappa}{1-\nu}$. 따라서 $G_\kappa(\nu)\ge G_\kappa(w)$이고, $\kappa(y)\to1$, $\eta(y)\to0$이므로 큰 $y$에서 $(1-\eta)G_{\kappa}(w)\ge(1-\eps)G_1(w)=(1-\eps)G(w)$.

\emph{경우 B: $\nu<\nu_*(y)$.} $c=1<b_j$에는 한 항이 필요하므로 $K_j\ge1$. 큰 $y$에서 $\nu\le1/2$이고 $s_j=\log\frac1{1-\nu}\le\frac{\nu}{1-\nu}\le2\nu<2\exp(-3(\log\log y)^2)$이므로 $G(w)s_j\log y\le2G(w)\log y\,\exp(-3(\log\log y)^2)\le1\le K_j$.
\end{proof}

\begin{theorem}[기하적 블록 방식의 $(\log L)^2$ 장벽]\label{STRUCT:thm-barrier}
\PROVED\ (정리~\ref{STRUCT:ext-pnt} 사용.) $0<w<1$, $Y_0\ge1$, $0<\eps<1$을 고정하고
\[
c(w):=\frac{G(w)}2=\frac1{2\,(1+\log\frac1w)\,\log\frac1{1-w}}
\]
라 하자. $L_0=L_0(w,Y_0,\eps)$가 있어, $L\ge L_0$이고 $\Lam_L$에 대한 $\Lam$형 블록 방식이 ``$y_j\ge Y_0$인 모든 $j\ge2$에서 $\nu_j\le w$''를 만족하면
\[
\sum_{j=1}^tK_j\ \ge\ H(y_1)+c(w)\bigl[(1-\eps)(\log L)^2-(\log y_1)^2\bigr].
\]
특히 $y_1\le L^{\vartheta}$이면 $\sum_jK_j\ge(1-\eps-\vartheta^2)c(w)(\log L)^2$.
\end{theorem}
\begin{proof}
$Y_1\ge Y_0$를 보조정리~\ref{STRUCT:lem-perblock} ($\eps/2$로)의 상수이면서 $y\ge Y_1$에서 $\psi(y)\le2y$가 되도록 잡는다. $K_1\ge H(y_1)$ (따름정리~\ref{STRUCT:cor-selfsim}). $t=1$이면 $y_1=L$이고 우변 괄호는 음수이므로 자명하다. $t\ge2$이면 $y_t=L\ge Y_1$이므로 $j_*:=\min\{j\ge2:y_j\ge Y_1\}$가 존재한다.
$j\ge j_*$에서 $y_j\ge Y_1$이고 $\nu_j\le w$이므로 보조정리~\ref{STRUCT:lem-perblock}로 $K_j\ge(1-\eps/2)G(w)s_j\log y_j$. $\tau\le\log\psi(y_j)$이면 $\tau\le\log(2y_j)$, 즉 $\log y_j\ge \tau-\log2$이므로
$s_j\log y_j\ge\int_{\log\psi(y_{j-1})}^{\log\psi(y_j)}(\tau-\log2)\,d\tau$.
$j\ge j_*$에 대해 합하면 적분구간이 이어져
\[
\sum_{j\ge j_*}K_j\ge(1-\tfrac\eps2)G(w)\int_{\log\psi(y_{j_*-1})}^{\log\psi(L)}(\tau-\log2)\,d\tau .
\]
$j_*=2$이면 $y_{j_*-1}=y_1$, 아니면 $y_{j_*-1}<Y_1$이므로 $Z:=\max(y_1,Y_1)$에 대해 $\psi(y_{j_*-1})\le\psi(Z)$이고, $\log\psi(Z)-\log2\le\log Z$, 또 큰 $L$에서 $\log\psi(L)-\log2\ge\log L-1$. 따라서 적분 $\ge\frac12[(\log L-1)^2-(\log Z)^2]\ge\frac12[(\log L)^2-2\log L-(\log y_1)^2-(\log Y_1)^2]$.
그러므로 $\sum_jK_j\ge H(y_1)+(1-\frac\eps2)c(w)[(\log L)^2-(\log y_1)^2]-c(w)[2\log L+(\log Y_1)^2]$이고, $L$이 크면 마지막 괄호 $\le\frac\eps2(\log L)^2$이므로 주장이 따른다.
\end{proof}

\begin{corollary}[DESIGN의 기하적 방식]\label{STRUCT:cor-geometric}
\PROVED\ $0<w'<1$을 고정하고 $\Lam$형 방식이 모든 $j$에서 $\log b_j\le w'y_j$를 만족한다고 하자. 그러면 모든 $\eps>0$에 대해 $L\ge L_0(w',\eps)$에서 $\sum_jK_j\ge(c(w')-\eps)(\log L)^2$.
수치: $c(1/4)=0.7283$, $c(1/2)=0.4260$, $c(3/4)=0.2801$, $c(0.9)=0.1964$; $w\to0$에서 $c(w)\sim1/(2w\log(1/w))$, $w\to1$에서 $c(w)\sim1/(2\log\frac1{1-w})\to0$.
\end{corollary}
\begin{proof}
$\psi(y)/y\to1$이므로 $Y(w')$가 있어 $y>Y(w')$이면 $\psi(y)>w'y$. $\log b_1=\psi(y_1)\le w'y_1$이므로 $y_1\le Y(w')$. $c$는 연속이므로 $w\in(w',1)$을 $c(w)\ge c(w')-\eps/2$로 잡는다. $y_j\ge Y_0(w,w')$이면 $\nu_j\le w'y_j/\psi(y_j)\le w$. 정리~\ref{STRUCT:thm-barrier}를 $\eps'$로 적용하면 $\sum K_j\ge c(w)[(1-\eps')(\log L)^2-(\log Y(w'))^2]\ge(c(w')-\eps)(\log L)^2$ (큰 $L$).
\end{proof}

\begin{remark}[수치: 최적 $\Lam$형 방식의 하한]\label{STRUCT:rem-num-block}
\NUMERIC\ \texttt{code/block\_opt.py} (출력 \texttt{out/block\_opt.out}): 소수거듭제곱 끝점을 갖는 모든 $\Lam$형 사슬에 대해 명제~\ref{STRUCT:prop-block-count}의 $\binom{D_j+K_j}{K_j}\ge b_j$ ($D_j$는 내림 반올림한 로그 히스토그램으로 위에서 누른 값)와 $K_1\ge H(y_1)$을 써서 $\sum K_j$의 최솟값을 동적계획으로 계산했다. $y_1\le22$, $\nu_j\le w$ ($j\ge2$)일 때 ($L$: $w=1/4,1/2,3/4,0.9$의 하한): $L=100$: $20,13,11,9$; $L=1000$: $40,26,21,17$; $L=5000$: $58,38,30,24$. 모두 $c(w)[(\log L)^2-(\log22)^2]$ ($L=5000$에서 $45.9,26.8,17.6,12.4$) 이상이고 $(\log L)^2$에 비례해 자란다. $y_1$을 제한하지 않으면 한 블록($y_1=L$, 엔트로피 하한 $5,\dots,11$)이 최소이다: 장벽은 첫 블록이 작은 방식에 관한 것이다.
\end{remark}

\begin{proposition}[장벽을 피하는 유일한 길: 각 척도에서의 SP1]\label{STRUCT:prop-escape}
\PROVED\ (a) $\Lam$형 방식에서 $K_1\ge H(y_1)$, $K_j\ge T(y_j,y_{j-1})$이고 $T(y,y')\le H(y)\le H(y')+T(y,y')$. 따라서 블록 $j$의 단계성질은 ``$\Lam_{y_j}$ 척도의 SP1을 인자 $\Lam_{y_{j-1}}=\Lam_{y_j}^{1-\nu_j}$만큼 뺀 것''이다: $H(y_j)\le H(y_{j-1})+K_j$.
(b) (한 블록) $t=1$이면 방식은 $H(L)\le K_1$ 자체, 곧 SP1이다. 정리~\ref{STRUCT:thm-barrier}에서 $\sum K_j\le A\log L$이 되려면 $(\log L)^2-(\log y_1)^2\le(A/c(w)+o(1))\log L$, 즉 $y_1\ge L\,\e^{-A/c(w)-o(1)}$이어야 하고 그때 $H(y_1)\le A\log L$, 곧 척도 $\asymp L$의 SP1이 필요하다.
(c) (초지수 블록) $0<\lambda<1$을 고정하면 다음은 동치이다:
(1) SP1, 곧 $\exists C\ \forall y\ge2:\ H(y)\le C\log y$;
(2) $\exists C',y_0\ \forall y\ge y_0:\ T(y,\lfloor y^\lambda\rfloor)\le C'\log y$;
(3) $\exists A$: 모든 $L$에서 사슬 $y_t=L$, $y_{j-1}=\lfloor y_j^\lambda\rfloor$ ($y_1<y_0$까지; $\nu_j\to1$)의 최적 비용 $H(y_1)+\sum_{j\ge2}T(y_j,y_{j-1})$이 $\le A\log L+A$.
즉 초지수 블록이 $O(\log L)$ 비용을 가지는 것은 각 블록이 자기 척도에서 SP1인 것과 같다.
(d) counting은 초지수 블록을 막지 않는다: $T(y,\lfloor y^\lambda\rfloor)\ge(1/\log2-o(1))\log y$이고, 이 counting 하한들의 사슬 합은 $\sum_j(1/\log 2)\log y_j\le\frac{\log L}{(1-\lambda)\log2}$ 이하이다.
\end{proposition}
\begin{proof}
(a) 따름정리~\ref{STRUCT:cor-selfsim}. (b) 앞 부분은 정의. $\sum K_j\le A\log L$이라 하자. 정리~\ref{STRUCT:thm-barrier} ($\eps=1/2$)에서 $(\log y_1)^2\ge\frac12(\log L)^2-\frac{A}{c(w)}\log L$이므로 $y_1\to\infty$. 따라서 임의의 고정 $\eps>0$에 대해 큰 $L$에서 $y_1\ge Y_1(w,\eps)$이고 모든 $j\ge2$가 보조정리~\ref{STRUCT:lem-perblock}의 조건을 만족한다. 정리~\ref{STRUCT:thm-barrier}의 증명을 $j_*=2$로 반복하면
$\sum_{j\ge2}K_j\ge(1-\eps)G(w)\int_{\log\psi(y_1)}^{\log\psi(L)}(\tau-\log2)d\tau=(1-\eps)c(w)\log\frac{\psi(L)}{\psi(y_1)}\bigl(\log\psi(L)+\log\psi(y_1)-2\log2\bigr)\ge(1-2\eps)c(w)\log\frac{\psi(L)}{\psi(y_1)}\log L$.
그러므로 $\log(\psi(L)/\psi(y_1))\le A/((1-2\eps)c(w))$이고 $\psi(y)\sim y$에서 $y_1\ge L\e^{-A/c(w)-o(1)}$; 또 $H(y_1)\le K_1\le A\log L$.
(c) (1)$\Rightarrow$(2): $T\le H$. (2)$\Rightarrow$(3): $\log y_j\le\lambda^{t-j}\log L$이므로 $\sum_{j\ge2}T(y_j,y_{j-1})\le C'\sum_{i\ge0}\lambda^i\log L=\frac{C'}{1-\lambda}\log L$, 그리고 $H(y_1)\le\max_{z<y_0}H(z)$. (3)$\Rightarrow$(1): 따름정리~\ref{STRUCT:cor-selfsim}로 $H(L)\le H(y_1)+\sum_{j\ge2}T(y_j,y_{j-1})\le A\log L+A$.
(d) $c<\Lam_y/\Lam_{\lfloor y^\lambda\rfloor}$를 덮으려면 $(1+\tau(\Lam_y))^K\ge\Lam_y/\Lam_{\lfloor y^\lambda\rfloor}=\e^{(1-o(1))y}$이고 $\log\tau(\Lam_y)\le\pi(y)\log2+\pi(\sqrt y)\log(1+\log_2y)=(\log2+o(1))y/\log y$ (PNT). 합은 $\log y_j\le\lambda^{t-j}\log L$에서.
\end{proof}

\begin{remark}[일반 사슬에는 이 장벽이 없다]\label{STRUCT:rem-general-chain}
\PROVED\ ($\Lam$형이 아닌 사슬) $N_j$가 $\Lam_L$의 임의의 약수이면 명제~\ref{STRUCT:prop-block-count}와 보조정리~\ref{STRUCT:lem-divcount} ($y=L$)는 $K_j\ge(1-o(1))\log L/(1+\log(\psi(L)/\log b_j))$만 준다. 같은 크기의 $t$개 블록이면 합은 $\ge(1-o(1))t\log L/(1+\log t)$이고 $t=2$에서 $O(\log L)$이다. 즉 $(\log L)^2$ 장벽은 블록 $j$가 더 작은 $\Lam_{y_j}$의 약수만 쓰는 최소공배수 사슬의 성질이다(풀 소수 부분집합으로 만든 사슬은 이 counting에 막히지 않는다; 그러한 사슬의 단계성질은 \OPEN).
\end{remark}

\begin{hypothesis}[$\mathrm{BC}^*(w,C)$]\label{STRUCT:hyp-BC}
$0<w<1$, $C>0$. $y_*$가 있어 모든 $y\ge y_*$와 모든 정수 $0\le c<\Lam_y^{\,w}=\e^{w\psi(y)}$에 대해 $g_y(c)\le C\log y$. \OPEN
\end{hypothesis}

\begin{theorem}[$\mathrm{BC}^*\Rightarrow(\log L)^2$]\label{STRUCT:thm-BC}
\COND{$\mathrm{BC}^*(w,C)$}\ $\displaystyle\limsup_{L\to\infty}\frac{H(L)}{(\log L)^2}\le\frac{C}{2\log\frac1{1-w}}$.
\end{theorem}
\begin{proof}
$0<\eps<w$를 고정한다. $Y\ge\max(y_*,3)$를 $y\ge Y$에서 $\log y\le\eps\psi(y)$, $\psi(y)\ge y/2$가 되도록 잡는다(PNT).
$z_0=L$로 두고 $z_i\ge Y$인 동안 $z_{i+1}:=\min\{z\in\Z_{\ge1}:\psi(z)\ge(1-w)\psi(z_i)\}$로 정한다. 그러면 $\psi(z_{i+1})\ge(1-w)\psi(z_i)$이고 $\psi(z_{i+1}-1)<(1-w)\psi(z_i)$, 또 $\psi(z)-\psi(z-1)\le\log z$이므로
\[
\psi(z_{i+1})<(1-w)\psi(z_i)+\log z_i\le(1-w+\eps)\psi(z_i)<\psi(z_i),
\]
따라서 $z_{i+1}<z_i$이고 $\Lam_{z_{i+1}}<\Lam_{z_i}$. $z_m<Y$가 되는 첫 $m$에서 멈추고 $\Lam$형 사슬 $1\mid\Lam_{z_m}\mid\Lam_{z_{m-1}}\mid\dots\mid\Lam_{z_0}=\Lam_L$을 쓴다.
첫 블록의 예산은 $H(z_m)\le H_Y:=\max_{z<Y}H(z)<\infty$ (따름정리~\ref{STRUCT:cor-practical}). 블록 $(z_{i+1},z_i]$ ($i<m$)에서는 $z_i\ge Y\ge y_*$이고 $b=\Lam_{z_i}/\Lam_{z_{i+1}}\le\e^{w\psi(z_i)}$이므로 모든 $c<b$는 $\mathrm{BC}^*$로 $\Lam_{z_i}$의 약수 $\le C\log z_i$개의 합이다. 정리~\ref{STRUCT:thm-block}로
\[
H(L)\le H_Y+C\sum_{i<m}\log z_i .
\]
$i<m$에서 $\psi(z_i)\le(1-w+\eps)^i\psi(L)$, $\log z_i\le\log(2\psi(z_i))\le\Lambda^\dagger-i\beta$ ($\Lambda^\dagger:=\log(2\psi(L))$, $\beta:=\log\frac1{1-w+\eps}$), 그리고 $\log z_i>0$. 따라서 $\sum_{i<m}\log z_i\le\sum_{0\le i\le\Lambda^\dagger/\beta}(\Lambda^\dagger-i\beta)\le\Lambda^\dagger+\frac{(\Lambda^\dagger)^2}{2\beta}$. $\Lambda^\dagger=\log L+O(1)$이므로 $\limsup H(L)/(\log L)^2\le C/(2\beta)$. $\eps\to0$.
\end{proof}

\begin{corollary}[$N(b)$와 Erdős \#18]\label{STRUCT:cor-BC-conseq}
\COND{$\mathrm{BC}^*(w,C)$}\ (i) $N(b)\le\bigl(\frac{C}{\log(1/(1-w))}+o(1)\bigr)(\log\log b)^2$.
(ii) $m=\Lam_L$은 모두 실용수이고 $h(m)=H(L)\le\bigl(\frac{C}{2\log(1/(1-w))}+o(1)\bigr)(\log\log m)^2$. 따라서 Erdős \#18의 첫 질문 (``$h(m)<(\log\log m)^{O(1)}$인 실용수 $m$이 무한히 많은가'', erdosproblems.com/18, snippet-verified, \$250)은 지수 $2$, 명시적 수열 $m=\Lam_L$, 모든 큰 $L$에 대해 긍정된다.
\end{corollary}
\begin{proof}
(i) 전달정리 (p1p4 정리 8.5, 따름정리 8.6, \PROVED): $N(b)\le2H(L_b)$, $L_b\le\max(7,\lceil\log_2b\rceil)$, $\log L_b\le\log\log b+O(1)$. (ii) 실용성은 따름정리~\ref{STRUCT:cor-practical}; $1\le n<m$을 합하는 약수는 자동으로 진약수이므로 $h(\Lam_L)=H(L)$; $\log\log\Lam_L=\log\psi(L)=\log L+o(1)$.
\end{proof}

\begin{remark}[$\mathrm{BC}^*$의 위치와 상수의 일치]\label{STRUCT:rem-BC}
\PROVED\ (1) SP1 ($H(y)\le C\log y$)은 모든 $w$에 대해 $\mathrm{BC}^*(w,C)$를 준다. $\mathrm{BC}^*$는 $(\log L)^2$만 주므로 형식적으로 더 약한 가설이지만, 길이 $\e^{w\psi(y)}$의 구간을 $O(\log y)$개의 약수로 덮으라는 지수적으로 세밀한 덮개 명제이며 \OPEN 이다 (알려진 무조건적 결과, 곧 작은 잔차 정리는 $\e^{O((\log y)^2)}$까지만 덮는다).
(2) counting (보조정리~\ref{STRUCT:lem-perblock}의 경우 A를 $\nu=w$에 적용): $\mathrm{BC}^*(w,C)$이면 $C\ge(1-o(1))/(1+\log\frac1w)$. 이 최소 $C$를 정리~\ref{STRUCT:thm-BC}에 넣으면 상수는 $\frac{1}{2(1+\log\frac1w)\log\frac1{1-w}}=c(w)$로 정리~\ref{STRUCT:thm-barrier}의 장벽 상수와 정확히 같다. 즉 기하적 방식 안에서 $(\log L)^2$와 그 상수는 counting 수준에서 최적이다.
\end{remark}
